\chapter{Classes with a non-closed background}
\label{chap:nonclosed-bclasses}

Removing a divisorial current leaves the exterior derivative of the background unchanged. The residual current therefore belongs to an affine space modeled on Bott--Chern cohomology, even when the background is not closed. Prescribing this derivative makes the classes compatible with pushforward, while the exceptional divisor of \cref{chap:divisorial-models} gives their change under pullback. Analytic singularities yield nef Cartier realizations; regularization approximates general singularities by such realizations with a controlled change of derivative.

Throughout the chapter, $Y$ is a connected compact complex manifold of dimension $n\geq1$, $\theta$ is a smooth real $(1,1)$-form on $Y$, and $\pi:Z\to Y$ denotes a smooth compact modification.

\section{An affine space on each model}\label{sec:affine-model-space}

Two currents with exterior derivative $\pi^*\mathrm{d}\theta$ differ by a closed current. Quotienting out global potential terms makes their difference a Bott--Chern class. Write $\mathcal D'(Z,\mathbb R)$ for real distributions and set
\[
 \mathscr A_Z(\theta)=
 \frac{\{\pi^*\theta+C:C\text{ is a closed real }(1,1)\text{-current on }Z\}}
      {\ddc\mathcal D'(Z,\mathbb R)}.
\]
Write $\{T\}$ for the class of $T$. By \cref{lem:bc-one-one-current-description}, every class has a smooth representative, and translation by Bott--Chern classes gives
\[
 \mathscr A_Z(\theta)=\{\pi^*\theta\}+H^{1,1}_{\mathrm{BC}}(Z,\mathbb R).
\]
The smooth forms representing elements of $\mathscr A_Z(\theta)$ are exactly those satisfying $\mathrm{d}\eta=\pi^*\mathrm{d}\theta$.

\Needspace{11\baselineskip}
\begin{proposition}\label{prop:fixed-background-class-maps}
If $g:Z'\to Z$ is a compact modification over $Y$, pushforward and pullback induce maps
\[
 g_*:\mathscr A_{Z'}(\theta)\to\mathscr A_Z(\theta),
 \qquad g^*:\mathscr A_Z(\theta)\to\mathscr A_{Z'}(\theta),
 \qquad g_*g^*=\mathrm{id}.
\]
Pushforward preserves classes with positive current representatives.
\end{proposition}
\begin{proof}
Write $\pi'=\pi\circ g$. For a representative $\pi'^*\theta+C'$, degree-one change of variables gives
\[
 g_*(\pi'^*\theta+C')=\pi^*\theta+g_*C'.
\]
The second term is closed. Changing the representative by $\ddc U$ changes its pushforward by $\ddc(g_*U)$, so the map is well-defined. Positivity follows from positivity of proper pushforward.

For pullback choose a smooth representative $\eta=\pi^*\theta+\beta$, where $\beta$ is smooth and closed. Its pullback has the required form on $Z'$. If two such representatives are equivalent, their difference is a smooth form $\ddc U$; the injectivity part of \cref{lem:bc-one-one-current-description} makes $U$ smooth. Their pullbacks consequently differ by $\ddc(U\circ g)$, proving independence. Finally $g_*g^*\eta=\eta$ for smooth representatives proves the identity on classes.
\end{proof}

The prescribed derivative is essential to this pushforward construction: $g_*\pi'^*\mathrm{d}\theta=\pi^*\mathrm{d}\theta$ is smooth, whereas pushing forward an arbitrary smooth exterior derivative can produce a singular current.

A compatible family with entries in $\mathscr A_Z(\theta)$ will be called a \emph{b-divisor class over the background $\theta$}. Compatibility always means the displayed pushforward compatibility. A family is pseudoeffective if every entry has a positive representative. A Cartier realization on $Z$ means that every entry on a model dominating $Z$ is the pullback of one class in $\mathscr A_Z(\theta)$.

An affine class is called nef if it admits smooth representatives $\eta_\epsilon\geq-\epsilon\gamma$ for every $\epsilon>0$, where $\gamma$ is a Hermitian metric on that model. Metric comparison makes the definition independent of $\gamma$. A Cartier realization is nef when its realizing affine class is nef.

\section{The class associated with trace data}\label{sec:affine-trace-class}

The trace currents have the required pushforward compatibility. Subtracting them from the background gives affine classes, and the common trace majorant supplies positive representatives on all models.

Suppose $Y$ is a smooth submanifold of a connected compact complex manifold $X$. Let $\Theta$ be a smooth real $(1,1)$-form on $X$ with $\Theta|_Y=\theta$, and let $u\in\PSH(X,\Theta)$ satisfy $\nu(u,Y)=0$. Use the trace currents $S_Z(u;Y)$ of \cref{sec:divisorial-model-bound}.

\Needspace{13\baselineskip}
\begin{theorem}\label{thm:nonclosed-trace-bclass}
For the trace data just fixed, the family defined by
\[
 \mathbb B_Z(u;Y,\theta)
 =\{\pi^*\theta-S_Z(u;Y)\}\in\mathscr A_Z(\theta)
\]
is a pseudoeffective b-divisor class over $\theta$. It is unchanged by ambient multiplier-ideal equivalence and by a smooth $\ddc$ change of the restricted background, after identifying the resulting affine spaces.
\end{theorem}
\begin{proof}
The currents $S_Z$ are closed, so the displayed classes belong to the indicated affine spaces. Their compatibility follows from \cref{thm:trace-divisorial-model-compatibility} and \cref{prop:fixed-background-class-maps}. \Cref{prop:original-background-trace-majorant} supplies one $v\in\PSH(Y,\theta)$ with
\[
 \pi^*\theta+\ddc(v\circ\pi)-S_Z(u;Y)\geq0
\]
on every model, proving pseudoeffectivity. The currents $S_Z$ depend only on the ambient multiplier ideals. For a smooth real function $f$ on $Y$, a change $\theta\mapsto\theta+\ddc f$ adds $\ddc(f\circ\pi)$ to every representative, which is zero in the quotient. This proves both invariances.
\end{proof}

\begin{corollary}\label{cor:affine-trace-exceptional-difference}
For the trace classes in \cref{thm:nonclosed-trace-bclass} and a compact modification $g:Z'\to Z$ over $Y$, let $E_g$ be the effective exceptional divisor in \cref{thm:trace-exceptional-difference}. Then
\[
 \mathbb B_{Z'}(u;Y,\theta)
       =g^*\mathbb B_Z(u;Y,\theta)-\{[E_g]\}_{\mathrm{BC}}.
\]
Here the Bott--Chern class on the right acts by translation on the affine space $\mathscr A_{Z'}(\theta)$. No closedness of $\theta$ is assumed.
\end{corollary}
\begin{proof}
Let $\pi:Z\to Y$ be the structure map. Abbreviate $S_Z=S_Z(u;Y)$ and write $S_Z=B_Z+\ddc v_Z$ with $B_Z$ smooth and closed and $v_Z$ quasi-plurisubharmonic. The affine class $\mathbb B_Z$ has smooth representative $\pi^*\theta-B_Z$. Its pullback therefore equals the class of $\pi'^*\theta-g^*S_Z$, where $\pi'=\pi\circ g$; the two representatives differ by $\ddc(v_Z\circ g)$. Substituting $S_{Z'}=g^*S_Z+[E_g]$ gives the displayed identity. The exterior derivative of both representatives is $\pi'^*\mathrm{d}\theta$, since the subtracted currents are closed. Thus the subtraction takes place within the stated affine space.
\end{proof}

The exceptional-divisor cocycle describes these translations under further modifications. The trace measures of their analytic-ideal approximants converge in total variation by \cref{thm:trace-exceptional-difference}. This is linear compatibility of classes with a prescribed derivative; it does not remove the separate problem of defining a compatible nonlinear volume.


When $\theta$ is closed, this family is the Bott--Chern family of \cref{prop:trace-bott-chern-bclass}.

For a potential $u\in\PSH(Y,\theta)$ on the base itself, let $E$ range over prime divisors on $Z$, write $[E]$ for the integration current and $\nu(u\circ\pi,E)$ for the generic Lelong number, and set
\[
 N_Z(u)=\sum_E\nu(u\circ\pi,E)[E],
 \qquad \mathbb D_\theta(u)_Z=\{\pi^*\theta-N_Z(u)\}.
\]
By \cref{lem:self-trace-is-siu}, $N_Z(u)=S_Z(u;Y)$. The preceding theorem makes $\mathbb D_\theta(u)$ a pseudoeffective family. Its positive representative can also be taken directly as $\pi^*\theta+\ddc(u\circ\pi)-N_Z(u)$ by \cref{lem:nonclosed-divisorial-removal} and passage to increasing finite sums.

\section{Analytic singularities and nef realization}

Resolving an analytic singularity leaves a bounded residual potential. Almost positive regularization makes its affine class nef, and boundedness prevents new divisorial terms on higher models.

\begin{proposition}\label{prop:nonclosed-analytic-nef-cartier}
If $u\in\PSH(Y,\theta)$ has analytic singularities, then $\mathbb D_\theta(u)$ has a nef Cartier realization on a smooth compact modification.
\end{proposition}
\begin{proof}
Resolve the analytic singularity ideal as in the proof of \cref{thm:analytic-current-nef-cartier}. On one model $\pi:Z\to Y$ we have
\[
 \pi^*\theta+\ddc(u\circ\pi)=[D]+R,
 \qquad R=\beta+\ddc v\geq0,
\]
where $D$ is a finite effective rational divisor, $\beta$ is smooth, $v$ is globally bounded after choosing divisor metrics, and $\mathrm{d}\beta=\pi^*\mathrm{d}\theta$. Positivity follows from divisorial removal.

Fix a Hermitian form $\gamma$ on $Z$ and apply \cref{prop:bounded-almost-regularization} to $v$ in the smooth background $\beta$. It gives smooth $v_j$ with $\beta+\ddc v_j\geq-\epsilon_j\gamma$, $\epsilon_j\downarrow0$. These forms represent the same affine class and have exterior derivative $\pi^*\mathrm{d}\theta$, proving nefness.

On every further modification, the pullback of the bounded residual potential remains bounded and has zero Lelong numbers. The divisorial part is precisely the pullback of $D$. Thus the residual affine class on that model is the pullback of the class of $R$, proving the Cartier realization.
\end{proof}

\section{Approximation while the derivative varies}

Analytic regularization retains a small curvature error. Absorbing it into the background allows the preceding proposition to apply. On each fixed model, convergence of the divisorial coefficients then gives convergence of residual-current representatives, with an explicit error in their exterior derivatives.

\Needspace{27\baselineskip}
\begin{theorem}\label{thm:varying-background-nef-approximation}
Let $u\in\PSH(Y,\theta)$ and fix a Hermitian form $\omega$ on $Y$. There are analytic potentials $u_m$ and positive numbers $\epsilon_m\downarrow0$, indexed by sufficiently large positive integers $m$, such that
\[
 \theta_m=\theta+2\epsilon_m\omega,
 \qquad \theta_m+\ddc u_m\geq\epsilon_m\omega.
\]
For each $m$, the family $\mathbb D_{\theta_m}(u_m)$ has a nef Cartier realization. On every fixed compact model $\pi:Z\to Y$, its class has the current representative
\[
 C_{m,Z}=\pi^*\theta+2\epsilon_m\pi^*\omega-D_{m,Z},
 \qquad D_{m,Z}=\frac1m\sum_E\operatorname{ord}_E\mathcal I(mu)[E].
\]
Here $E$ ranges over prime divisors on $Z$, and the order is computed after pulling the ideal back to $Z$. These representatives satisfy
\[
 C_{m,Z}\longrightarrow\pi^*\theta-N_Z(u)
 \quad\hbox{weakly as currents},
\]
and their exterior derivatives converge smoothly to $\pi^*\mathrm{d}\theta$.
\end{theorem}
\begin{proof}
Take $u_m$ and $\epsilon_m$ from \cref{cor:nonclosed-analytic-approximation}, enlarging the errors to make them positive if necessary. Adding $2\epsilon_m\omega$ changes the lower bound $-\epsilon_m\omega$ into $\epsilon_m\omega$, proving the first assertion. Each $u_m$ has analytic type $(1/m,\mathcal I(mu))$, so \cref{prop:nonclosed-analytic-nef-cartier} applies to its positive current in background $\theta_m$.

On a fixed model the generic divisorial Lelong numbers of $u_m\circ\pi$ are exactly $m^{-1}\operatorname{ord}_E\mathcal I(mu)$. Dropping the $\ddc$ of that potential from the residual-current representative leaves precisely $C_{m,Z}$ in its affine class. The self-trace identity and \cref{thm:trace-divisorial-current-budget} give $D_{m,Z}\to N_Z(u)$ in trace mass, hence weakly as currents. The smooth error $2\epsilon_m\pi^*\omega$ tends to zero in every derivative, proving current convergence. Since each $D_{m,Z}$ is closed,
\[
 \mathrm{d}C_{m,Z}=\pi^*\mathrm{d}\theta+2\epsilon_m\pi^*\mathrm{d}\omega,
\]
which proves the stated smooth convergence.
\end{proof}

The representatives converge as currents, but their classes lie in $\mathscr A_Z(\theta_m)$. If $\mathrm{d}\omega\neq0$, their derivatives differ from the prescribed derivative of $\mathscr A_Z(\theta)$ by $2\epsilon_m\pi^*\mathrm{d}\omega$. Thus even for closed $\theta$, this construction does not give approximation by closed nef Cartier classes. Approximation with exactly fixed derivative remains a separate problem.

\section{Divisorial classes and infinite mass}

The infinite-mass potential of \cref{thm:infinite-mass-witness} has trivial multiplier ideals at every weight. Its divisorial data on all modifications therefore agree with those of a smooth potential. This gives finite and infinite extended masses within one nef Cartier family.

Take the background to be a Hermitian form $\omega$ on $Y$, with upper mass $v_+(\omega)$ and extended measure $\mathrm{MA}^{\mathrm{ext}}_\omega$ as in \cref{chap:products}.

\Needspace{22\baselineskip}
\begin{proposition}\label{prop:affine-bclass-infinite-mass}
If $v_+(\omega)=+\infty$, there is $U\in\PSH(Y,\omega)$ such that
\[
 \mathbb D_\omega(U)=\mathbb D_\omega(0),\qquad
 \int_Y\mathrm{MA}^{\mathrm{ext}}_\omega(U)=+\infty,
 \qquad \int_Y\omega^n<+\infty.
\]
The common family has the nef Cartier realization $(Y,\{\omega\})$.
\end{proposition}
\begin{proof}
Take $U$ from \cref{thm:infinite-mass-witness}. Every positive-weight multiplier ideal is the unit ideal. Hence every divisorial order in $D_{m,Z}$ is zero on every model. The self-trace identity gives $N_Z(U)=0$, which is also the divisorial part of the constant potential. Both families are therefore the pullback classes of $\omega$. The form $\omega$ is smooth and positive, so its affine class is nef and realizes the family on the base. The infinite extended mass is part of the construction of $U$, whereas the smooth measure $\omega^n$ has finite mass by compactness.
\end{proof}

Intersection products for these affine families and their relation to finite damped measures remain to be constructed. The preceding example shows that such a construction cannot recover the raw mass of every potential from its residual family alone.
